% Prepared for Edward Baker, 22 October 2026.
% LLM assistance: gpt-6-astra (Codex), effort ultra, inherited parent configuration.
% Complete elementary concentration proof; no differentiated saddle result
% or unproved pair-correlation conjecture is used in these results.

\section{Complete-source theta clock limits and growing windows}
\label{pm:theta-section}
\label{th:section}

This section supplies a self-contained canonical application of the finite
matrix theorem.  The proof uses the entire positive theta source, an
explicit elementary concentration bound, and a real Fourier comparison.
It proves locations, simplicity, rank identification and matrix
convergence at large real heat time, including the growing windows
specified below. It does not differentiate a root
location remainder.  A separate, stronger differentiated saddle theorem
would be needed for pointwise velocity asymptotics; those asymptotics are
not used here.

\subsection{The source and its elementary bounds}
\label{pm:theta-source-subsection}

Use the fixed normalization
\begin{equation}
 H(t,z)=\int_0^\infty e^{tu^2}\Phi(u)\cos(zu)\,du,
 \qquad
 \Phi(u)=\sum_{n=1}^\infty
       \bigl(2\pi^2n^4e^{9u}-3\pi n^2e^{5u}\bigr)
                         e^{-\pi n^2e^{4u}}.
 \label{pm:theta-source}
\end{equation}
In particular the complete summands are positive for \(u\ge0\).
Writing
\begin{align}
 \Phi(u)&=2\pi^2e^{9u-\pi e^{4u}}\mathfrak b(u),
 \label{pm:theta-factorization}\\
 \mathfrak b(u)
  &=\sum_{n=1}^\infty
       \left(n^4-\frac{3n^2}{2\pi}e^{-4u}\right)
                       e^{-\pi(n^2-1)e^{4u}},
 \label{pm:theta-b-definition}
\end{align}
we have the global bounds
\begin{equation}
 0<\mathfrak b_-:=1-\frac3{2\pi}
 \le\mathfrak b(u)\le
 \mathfrak b_+:=\frac1{1-16e^{-3\pi}}<\infty.
 \label{pm:theta-b-bounds}
\end{equation}
Indeed the \(n=1\) term supplies the lower bound, and every term is
positive.  For the upper bound discard only the negative summand in this
upper estimate and use
\(n\le2^{n-1}\), \(n^2-1\ge3(n-1)\), and \(e^{4u}\ge1\).
The resulting series is bounded by the geometric series with ratio
\(16e^{-3\pi}<1\).  No term is discarded from the field in
\eqref{pm:theta-source} itself.

The upper bound in \eqref{pm:theta-b-bounds} implies convergence of
\eqref{pm:theta-source} and all its parameter derivatives uniformly on
compact subsets of \(\mathbb C^2\).  In fact every such derivative is
bounded in absolute value by a constant times
\[
 \int_0^\infty (1+u)^j
             e^{A u^2+B u+9u-\pi e^{4u}}\,du<\infty
\]
for suitable finite \(A,B,j\).  Thus \(H\) is holomorphic in both
variables and satisfies \(H_t=-H_{zz}\) by differentiation under the
integral.  For real \(t\), its normalizing value \(H(t,0)\) is
strictly positive.

For \(L\ge2\), define
\begin{equation}
 t=t(L)=\frac{4\pi e^{4L}-9}{2L},\qquad
 \kappa=16\pi e^{4L}-2t=8tL-2t+36.
 \label{pm:theta-saddle-parameters}
\end{equation}
The function \(t(L)\) is strictly increasing, because
\[
 t'(L)=\frac{4\pi e^{4L}(4L-1)+9}{2L^2}>0.
\]
It maps \([2,\infty)\) onto \([t_2,\infty)\), where
\(t_2=(4\pi e^8-9)/4\).  We use its inverse to define \(L=L(t)\)
for every \(t\ge t_2\).  Set
\begin{equation}
 q_t(u)=tu^2+9u-\pi e^{4u}.
 \label{pm:theta-q-definition}
\end{equation}
Then \(q_t'(L)=0\) and \(q_t''(L)=-\kappa<0\).  The identities
\eqref{pm:theta-saddle-parameters} imply
\begin{equation}
 \kappa\ge36,\qquad
 \frac t\kappa\le\frac1{8L-2}\le\frac1{14}<\frac1{4e},
 \qquad \pi e^{4L}\le\frac\kappa8.
 \label{pm:theta-parameter-bounds}
\end{equation}
For the last inequality, writing \(A=\pi e^{4L}\), one has
\(t=(4A-9)/(2L)<A\) for \(L\ge2\), and hence
\(\kappa=16A-2t>14A\).

\subsection{An explicit first-moment concentration theorem}
\label{pm:theta-concentration-subsection}

Let \(U\) denote a nonnegative random variable with the complete
probability law
\begin{equation}
 d\mu_t(u)=\frac{e^{tu^2}\Phi(u)}{H(t,0)}\,du.
 \label{pm:theta-probability}
\end{equation}
This notation is only a convenient integral normalization; the zero
configuration itself is deterministic.  Put \(c_*=(4e)^{-1}\).

\begin{theorem}[Complete elementary concentration bound]
\label{pm:theta-concentration}
For every \(t\ge t_2\), with \(L,\kappa\) as in
\eqref{pm:theta-saddle-parameters},
\begin{equation}
 \mathbb E_t|U-L|\le M_1(t):=
 \frac{e\mathfrak b_+}{\mathfrak b_-}
 \left\{\frac1{c_*\sqrt\kappa}
       +\frac{L^2\sqrt\kappa}{2}
                   e^{-c_*\kappa/16}\right\}.
 \label{pm:theta-M1-bound}
\end{equation}
In particular
\begin{equation}
 M_1(t)=O(\kappa^{-1/2}),\qquad
 L\sim\tfrac14\log t,\qquad \kappa\sim8tL
 \quad(t\to\infty).
 \label{pm:theta-concentration-asymptotics}
\end{equation}
The bound includes every theta mode and the entire half-line of
integration.
\end{theorem}

\begin{proof}
Write \(A=\pi e^{4L}\), \(v=u-L\), and \(a=1/4\).
The saddle equation gives the exact identity
\begin{equation}
 q_t(L+v)-q_t(L)=tv^2-A(e^{4v}-1-4v).
 \label{pm:theta-exact-exponent}
\end{equation}
For \(v\ge0\), Taylor's formula with positive remainder gives
\(e^{4v}-1-4v\ge8v^2\), so
\begin{equation}
 q_t(L+v)-q_t(L)\le-\frac\kappa2v^2.
 \label{pm:theta-right-gaussian}
\end{equation}
For \(-a\le v\le0\), the integral second-order Taylor remainder
instead gives \(e^{4v}-1-4v\ge8e^{-1}v^2\).  Since
\(8A=(\kappa+2t)/2\),
\[
 t-8A/e=-\frac\kappa{2e}+t(1-e^{-1})
 \le-\frac\kappa{4e}=-c_*\kappa
\]
by \eqref{pm:theta-parameter-bounds}.  Combining both sides yields
\begin{equation}
 q_t(L+v)-q_t(L)\le-c_*\kappa v^2
            \qquad(-a\le v<\infty).
 \label{pm:theta-near-gaussian}
\end{equation}

We must also pay the whole far-left interval \(0\le u\le L-a\).
The second derivative \(q_t''(u)=2t-16\pi e^{4u}\) is strictly
decreasing.  Any local maximum of \(q_t\) must be on the descending
part of \(q_t'\), where its only zero is the already identified
point \(L\).  Thus \(q_t\) has no interior local maximum in
\([0,L-a]\); its maximum there is at an endpoint.  Moreover,
\begin{align*}
 q_t(L-a)
 &=t\left((L-\tfrac14)^2-\frac{L}{2e}\right)
       +9\left(L-\tfrac14-\frac1{4e}\right)>0,
\end{align*}
where both brackets are positive for \(L\ge2\).  Since
\(q_t(0)=-\pi<0\), the right endpoint is the larger one.
Equation~\eqref{pm:theta-near-gaussian} at \(v=-a\) consequently
gives the complete far-left bound
\begin{equation}
 q_t(u)-q_t(L)\le-\frac{c_*\kappa}{16}
                  \qquad(0\le u\le L-\tfrac14).
 \label{pm:theta-far-left}
\end{equation}

For a normalization lower bound, use
\(|v|\le1/(2\sqrt\kappa)\).  This interval lies in \([-a,a]\)
and in the physical half-line after translation by \(L\).
The elementary upper Taylor bound
\[
 0\le e^{4v}-1-4v\le8e^{4|v|}v^2
\]
and \(A\le\kappa/8\) show that
\[
 q_t(L+v)-q_t(L)\ge-\kappa e^{4|v|}v^2
             \ge-e/4>-1.
\]
Here \(4|v|\le2/\sqrt\kappa<1\).  Therefore
\begin{equation}
 Z_t:=\int_0^\infty e^{q_t(u)-q_t(L)}\mathfrak b(u)\,du
       \ge\frac{\mathfrak b_-e^{-1}}{\sqrt\kappa}.
 \label{pm:theta-normalization-lower}
\end{equation}

Using \eqref{pm:theta-near-gaussian},
\eqref{pm:theta-far-left}, and \(\mathfrak b\le\mathfrak b_+\),
the first absolute-moment numerator is at most
\begin{align*}
 &\int_0^\infty |u-L|e^{q_t(u)-q_t(L)}\mathfrak b(u)\,du\\
 &\quad\le\mathfrak b_+
 \left\{\int_{\mathbb R}|v|e^{-c_*\kappa v^2}\,dv
        +e^{-c_*\kappa/16}\int_{-L}^{-1/4}|v|\,dv\right\}\\
 &\quad\le\mathfrak b_+
       \left\{\frac1{c_*\kappa}
                  +\frac{L^2}{2}e^{-c_*\kappa/16}\right\}.
\end{align*}
Division by \eqref{pm:theta-normalization-lower} proves
\eqref{pm:theta-M1-bound}.

Finally,
\[
 t=\frac{2\pi e^{4L}}L-\frac9{2L},\qquad
 \kappa=4\pi e^{4L}(4-L^{-1})+9L^{-1}.
\]
These formulas prove \(L\sim(\log t)/4\) and
\(\kappa\sim8tL\).  They also give
\(\kappa\ge14\pi e^{4L}\) for \(L\ge2\), and therefore
\(L^2\kappa e^{-c_*\kappa/16}\to0\).  The second term in
\eqref{pm:theta-M1-bound} is consequently
\(o(\kappa^{-1/2})\), proving its stated order.
\end{proof}

\subsection{The Fourier transfer and the first finitely many roots}
\label{pm:theta-fourier-subsection}

The following simple transfer isolates exactly which information from
concentration is used.  It does not assume real-rootedness of the entire
function.

\begin{lemma}[Real Fourier concentration transfer]
\label{pm:fourier-transfer}
Let \(U\ge0\) have a probability law with finite first moment, let
\(L>0\), and let \(m_1=\mathbb E|U-L|\).  For real \(z\) put
\(F(z)=\mathbb E\cos(zU/L)\).  Then for every \(Z>0\),
\begin{equation}
 \sup_{|z|\le Z}|F(z)-\cos z|\le\frac{Zm_1}{L},\qquad
 \sup_{|z|\le Z}|F'(z)+\sin z|\le\frac{(1+Z)m_1}{L}.
 \label{pm:fourier-C1-bounds}
\end{equation}
\end{lemma}

\begin{proof}
Write \(U/L=1+V/L\), with \(V=U-L\).  Since cosine is
\(1\)-Lipschitz on the real line,
\[
 |\cos(zU/L)-\cos z|\le |z||V|/L.
\]
The finite first moment permits differentiation under the expectation.
For the derivative use
\begin{align*}
 \left|\frac UL\sin(zU/L)-\sin z\right|
 &\le\frac{|V|}{L}|\sin(zU/L)|
           +|\sin(zU/L)-\sin z|\\
 &\le(1+|z|)|V|/L.
\end{align*}
Taking expectations proves both bounds.
\end{proof}

\begin{theorem}[Canonical fixed-rank clock and matrix limit]
\label{pm:theta-clock}
Fix \(K\ge1\) and \(0<\beta<\pi/4\), and set
\begin{equation}
 s_j=(j-\tfrac12)\pi,\qquad Z=s_K+\beta,\qquad
 \epsilon_0(t)=\frac{ZM_1(t)}L,\quad
 \epsilon_1(t)=\frac{(1+Z)M_1(t)}L.
 \label{pm:theta-finite-errors}
\end{equation}
Suppose \(t\ge t_2\) and
\begin{equation}
 \epsilon_0(t)<\sin\beta,\qquad
 \epsilon_1(t)<\cos\beta.
 \label{pm:theta-finite-tests}
\end{equation}
Then the first \(K\) positive real zeros of \(H(t,\cdot)\) exist,
are simple, and satisfy
\begin{equation}
 |Lx_j(t)-s_j|\le e(t):=\arcsin\epsilon_0(t)<\beta,
             \qquad1\le j\le K.
 \label{pm:theta-root-enclosure}
\end{equation}
All these finite tests hold for every sufficiently large \(t\), and
\begin{equation}
 e(t)=O_K((L\sqrt\kappa)^{-1}).
 \label{pm:theta-root-rate}
\end{equation}
No RH assumption is needed for these assertions about the indicated
late-time positive roots.

In particular, fix \(\ell\ge0\), \(m\ge2\), \(K=\ell+m\), and
put
\begin{equation}
 y_i=\frac L\pi x_{\ell+i}(t),\qquad
 c_i=\ell+i-\tfrac12,\qquad \delta(t)=e(t)/\pi.
 \label{pm:theta-clock-coordinates}
\end{equation}
The finite matrix bounds \eqref{pm:loewner-comparison} and
\eqref{pm:operator-comparison} hold with these data.  Consequently
\begin{align}
 \|\mathcal L(y)-\mathcal L_*\|_{\mathrm{op}}
    &=O_K((L\sqrt\kappa)^{-1}),
 \label{pm:theta-normalized-matrix-rate}\\
 \left\|\mathcal L(x_{\ell+1},\ldots,x_{\ell+m})
                -\frac{L^2}{\pi^2}\mathcal L_*\right\|_{\mathrm{op}}
    &=O_K(L/\sqrt\kappa).
 \label{pm:theta-physical-matrix-rate}
\end{align}
The second error tends to zero, while the leading physical matrix grows
like \(L^2\).
\end{theorem}

\begin{proof}
Normalize the actual field by
\[
 F_t(z)=\frac{H(t,z/L)}{H(t,0)}=\mathbb E_t\cos(zU/L).
\]
Theorem~\ref{pm:theta-concentration} and
Lemma~\ref{pm:fourier-transfer} give
\[
 |F_t(z)-\cos z|\le\epsilon_0,
 \qquad |F_t'(z)+\sin z|\le\epsilon_1
                     \quad(|z|\le Z).
\]
On the interval \([s_j-\beta,s_j+\beta]\), the endpoint cosine
values have opposite signs and absolute value \(\sin\beta\).
The first strict test in \eqref{pm:theta-finite-tests} preserves those
signs for \(F_t\).  On that interval the function \(\sin z\) has
one sign and absolute value at least \(\cos\beta\).  The second
strict test therefore gives a nonzero derivative with fixed sign for
\(F_t\).  There is exactly one simple zero in each interval.
On the complement of these intervals in \([0,Z]\), one has
\(|\cos z|\ge\sin\beta\), so the first test excludes every other
positive real zero there.  This proves the rank statement.

At the zero in the \(j\)-th interval,
\(|\cos z|=|\sin(z-s_j)|\le\epsilon_0\).  Since
\(|z-s_j|<\beta<\pi/2\), inversion of sine gives
\eqref{pm:theta-root-enclosure}.  The concentration theorem shows that
both \(\epsilon_0\) and \(\epsilon_1\) tend to zero, with size
\(O_K((L\sqrt\kappa)^{-1})\); hence the tests eventually hold and
so does \eqref{pm:theta-root-rate}.

For the selected block, \(|y_i-c_i|\le\delta<1/4\).  Apply
Theorem~\ref{pm:clock-matrix}.  Since \(m\) is fixed, its finite sum
in \eqref{pm:operator-comparison} is bounded uniformly for
\(0\le\delta<1/4\), proving
\eqref{pm:theta-normalized-matrix-rate}.  The physical rate follows
by multiplying that bound by \(L^2/\pi^2\), using
\eqref{pm:spatial-scaling}.
\end{proof}

The proof establishes simplicity on these large-time real root blocks.
It does not establish simplicity at time zero, nor does it identify a
hypothetical initial multiple root with one of these labels without
additional global continuation hypotheses.  The asymptotic statements
above fix the maximum rank \(K\).  The explicit finite tests also permit
the following uniform spatial consequence.

\begin{corollary}[Growing clock windows and dimension-independent matrix error]
\label{pm:theta-growing-clock}
Fix $0<\beta<\pi/4$ and let $K=K(t)\ge1$ be integer valued. At each
time define $Z,\epsilon_0,\epsilon_1$ by
\eqref{pm:theta-finite-errors} with this value of $K$. Whenever
\eqref{pm:theta-finite-tests} holds, all first $K$ positive real roots
are simple and obey \eqref{pm:theta-root-enclosure} simultaneously.
For every contiguous block of $m\ge2$ roots ending by rank $K$, set
$y_i=Lx_{\ell+i}/\pi$, $c_i=\ell+i-1/2$, and
$\delta=\pi^{-1}\arcsin\epsilon_0<1/4$. Then
\begin{align}
 \|\mathcal L(y)-\mathcal L(c)\|_{\mathrm{op}}
 &\le\frac{16\zeta(3)}{(1-2\delta)^3}\delta
 \le128\zeta(3)\delta.\label{pm:growing-matrix-bound}
\end{align}
In particular, the sufficient condition
\begin{equation}\label{pm:growing-window}
 K(t)=o\bigl(L(t)\sqrt{\kappa(t)}\bigr)
\end{equation}
makes the finite tests hold eventually and gives, uniformly over all
such blocks,
\begin{align}
 \max_{j\le K}|Lx_j-s_j|
   &=O\!\left(\frac K{L\sqrt\kappa}\right)=o(1),
       \label{pm:growing-root-rate}\\
 \|\mathcal L(y)-\mathcal L(c)\|_{\mathrm{op}}
   &=O\!\left(\frac K{L\sqrt\kappa}\right),
       \label{pm:growing-normalized-rate}\\
 \left\|\mathcal L(x_{\ell+1},\ldots,x_{\ell+m})
             -\frac{L^2}{\pi^2}\mathcal L(c)\right\|_{\mathrm{op}}
   &=O\!\left(\frac{KL}{\sqrt\kappa}\right).
       \label{pm:growing-physical-rate}
\end{align}
The constants do not depend on $K$, $m$ or the block's starting rank.
The stronger condition $K=o(\sqrt\kappa/L)$ is sufficient for the
absolute physical-matrix error to tend to zero.
\end{corollary}
\begin{proof}
The sign, monotonicity and complementary-interval arguments in
Theorem~\ref{pm:theta-clock} are finite-time implications with explicit
$K$ dependence, so apply to $K(t)$ at each time. For $d\ge1$,
$d-2\delta\ge(1-2\delta)d$, and therefore
\[
 \sum_{d=1}^{m-1}(d-2\delta)^{-3}
 \le(1-2\delta)^{-3}\sum_{d=1}^\infty d^{-3}
 =(1-2\delta)^{-3}\zeta(3).
\]
Insert this in \eqref{pm:operator-comparison}; $\delta<1/4$ makes
$(1-2\delta)^{-3}\le8$. This proves
\eqref{pm:growing-matrix-bound} without a population factor.
Since $Z\le\pi K$ and $1+Z\le(1+\pi)K$, the
$K$-independent bound $M_1=O(\kappa^{-1/2})$ proves the eventual
tests and the first two rates. Multiplication by $L^2/\pi^2$ gives
the physical rate. Finally $L\to\infty$, so the stronger condition
also implies \eqref{pm:growing-window}.
\end{proof}

At each time the compared matrices have the same dimension, which may
grow; no operator limit between unspecified varying spaces is asserted.
The corollary is spatial and requires no time regularity of $K(t)$.
Changing the population in time creates entry and exit terms. In
particular it does not extend the next fixed-population discriminant
identity by substituting $m(t)$, nor does it estimate the complete
centered exterior residual or yield a pair-correlation law.

\subsection{Cumulative exterior cancellation without velocity asymptotics}
\label{pm:theta-cancellation-subsection}

The location theorem also proves a useful dynamical consequence.  It is
weaker than a pointwise estimate for \((\log\Delta)'\), but requires
no differentiated saddle analysis and exposes the necessary exterior
payment in a fully analytic way.

\begin{corollary}[Physical energy growth and cumulative cancellation]
\label{pm:theta-cumulative-cancellation}
Fix the block from Theorem~\ref{pm:theta-clock} and choose \(T_0\)
large enough that its strict finite tests hold for every \(t\ge T_0\).
The root labels are real analytic on this time interval.  Define
\(\Delta(t)\), \(P_{\mathrm{pair}}(t)\), and
\(B_{\mathrm{ext}}(t)\) by
\eqref{pm:energy-definition} and
\eqref{pm:complete-discriminant-balance}, with the complete local
unit of the actual theta field.  Put
\begin{equation}
 A_m=\sum_{d=1}^{m-1}\frac{m-d}{d^2},\qquad
 C_m=m(m-1)\log\pi
       +2\sum_{d=1}^{m-1}(m-d)\log d.
 \label{pm:theta-pair-constants}
\end{equation}
Then
\begin{align}
 \log\Delta(t)&=-m(m-1)\log L(t)+C_m+o(1),
 \label{pm:theta-logdelta-asymptotic}\\
 P_{\mathrm{pair}}(t)&=\frac{8A_m}{\pi^2}L(t)^2
                         +O_K\!\left(\frac{L(t)}{\sqrt{\kappa(t)}}\right),
 \label{pm:theta-pair-asymptotic}\\
 \int_{T_0}^T P_{\mathrm{pair}}(t)\,dt
       &\sim\frac{8A_m}{\pi^2}T L(T)^2,
 \label{pm:theta-integrated-pair}\\
 \int_{T_0}^T B_{\mathrm{ext}}(t)\,dt
       &=-\int_{T_0}^T P_{\mathrm{pair}}(t)\,dt
             -m(m-1)\log L(T)+O(1).
 \label{pm:theta-integrated-exterior}
\end{align}
In particular \(E(t)=-\log\Delta(t)\to+\infty\), and
\begin{equation}
 \frac{\displaystyle\int_{T_0}^T B_{\mathrm{ext}}(t)\,dt}
      {\displaystyle\int_{T_0}^T P_{\mathrm{pair}}(t)\,dt}
                  \longrightarrow-1.
 \label{pm:theta-cancellation-ratio}
\end{equation}
These are cumulative statements.  They assert no sign for the pointwise
remainder obtained by differentiating
\eqref{pm:theta-logdelta-asymptotic}.
\end{corollary}

\begin{proof}
The field is holomorphic in both variables, and each selected zero is
simple.  The implicit-function theorem gives local analytic labels;
the unique rank identification in Theorem~\ref{pm:theta-clock} makes
these labels agree on overlapping intervals.

For \(i<j\), the root enclosure gives
\[
 x_{\ell+j}-x_{\ell+i}
  =\frac\pi L\left((j-i)+O_K((L\sqrt\kappa)^{-1})\right).
\]
There are finitely many pairs.  Taking logarithms and summing proves
\eqref{pm:theta-logdelta-asymptotic}, with \(C_m\) as displayed.
The inverse-square expansion, or equivalently the trace of
\eqref{pm:theta-physical-matrix-rate}, gives
\eqref{pm:theta-pair-asymptotic}.

The relation \(L(t)\sim(\log t)/4\) implies
\[
 \int_{T_0}^T L(t)^2\,dt\sim\frac1{16}T(\log T)^2
                         \sim T L(T)^2.
\]
To see this directly, fix any relative error, use the asymptotic
comparison with \((\log t)^2/16\) beyond a fixed threshold, and
integrate the upper and lower bounds; the finite initial interval is
negligible.  The elementary antiderivative
\(\int(\log t)^2dt=t((\log t)^2-2\log t+2)\)
then proves the assertion.  Applying the same comparison argument to
\(P_{\mathrm{pair}}(t)/L(t)^2\to8A_m/\pi^2\) proves
\eqref{pm:theta-integrated-pair}.

Finally the exact identity
\eqref{pm:complete-discriminant-balance} integrates to
\[
 \int_{T_0}^T B_{\mathrm{ext}}(t)\,dt
  =\log\Delta(T)-\log\Delta(T_0)
          -\int_{T_0}^T P_{\mathrm{pair}}(t)\,dt.
\]
Substitute \eqref{pm:theta-logdelta-asymptotic} to obtain
\eqref{pm:theta-integrated-exterior}.  The logarithmic remainder is
negligible relative to the positive quantity in
\eqref{pm:theta-integrated-pair}, proving
\eqref{pm:theta-cancellation-ratio}.  The same log-discriminant formula
proves \(E(t)\to+\infty\).
\end{proof}

The complete finite-cluster unit therefore cannot be omitted when the
inverse-square matrix is used dynamically.  Its cumulative contribution
cancels an order-\(T L(T)^2\) internal pair contribution, leaving an
order-\(\log L(T)\) change in the log discriminant.  This conclusion is
specific to the proved fixed-rank clock regime and the complete field.
It is neither an exclusion of other functionals nor a proof of
large-population random-matrix statistics.

\subsection{Limits of the statistical interpretation}
\label{pm:theta-scope-subsection}

For the fixed block, the normalized ordered pair differences tend to the
integers \(j-i\).  For example its empirical measure
\[
 \nu_{m,t}=\frac1{m(m-1)}\sum_{i\ne j}\delta_{y_j(t)-y_i(t)}
\]
converges weakly to
\[
 \frac1{m(m-1)}\sum_{i\ne j}\delta_{j-i}.
\]
This follows immediately by evaluating any bounded continuous test
function on the finitely many converging differences.  It is an atomic
clock limit.  It is not a claim about a smooth limiting pair-correlation
density, a random matrix ensemble, or a high-height limit at fixed heat
time.  Nor does fixed-rank convergence imply convergence uniformly over
all roots.  Those populations and limiting operations have to be
specified and estimated separately.

Finally, the full-source proof above exploits definite analytic
properties of the canonical source, but its conclusion alone is not a
simplicity discriminator at an earlier time.  The exact Hermite birth in
\eqref{pm:exact-hermite-birth} already exhibits simple late roots and a
multiple initial zero.  In the present theta regime the late physical gaps
shrink, while the normalized gaps converge to positive values.  Neither
feature supplies a vanishing complex-time boundary integral or pays an
unmeasured closure face.  Such a boundary consumer still needs the
complete field, its exterior unit and amplitude, its chosen topology,
and every term in its own transport identity.
